\documentclass{../zkdl-presentation-template}

\title[Криптографія]{\textbf{Практичне №5.Асиметрична криптографія}}

\date{7.10.2025}


\begin{document}
    \frame{
        \titlepage
    }
  
    \begin{frame}{Plan}
        \tableofcontents
    \end{frame}

\section{Теорія груп, кілець, полів} % Note all sections and subsections are automatically placed in your table of contents

%------------------------------------------------
\begin{frame}
    \frametitle{Теорія груп}
    \begin{block}{Означення}
        Групою $(G,*)$ називається алгебраїчна структура з однією операцією $*$, для якої виконується умови:
        \begin{enumerate}
            \item операція $*$ асоціативна: $\forall a,b,c\in G$
            \begin{equation*}
                a*(b*c)=(a*b)*c
            \end{equation*}
            \item існує нейтральний елемент $e\in G$ такий, що для будь-якого $a\in G$:
            \begin{equation*}
                e*a=a*e=a
            \end{equation*}
            \item для будь-якого $a\in G$ існує $a^{-1}\in G$ такий, що 
            \begin{equation*}
                a*a^{-1}=a^{-1}*a=e
            \end{equation*}
        \end{enumerate}
    \end{block}

    \vspace{-2mm}
    \begin{alertblock}{}
        Якщо виконується властивість 1) -- структура є \textbf{напівгрупою}, якщо виконується 1)+2) -- структура є \textbf{моноїдом}, якщо виконується 1)+2)+3)+комутативніть -- група є \textbf{абелевою}.
    \end{alertblock}
\end{frame}
%------------------------------------------------
\begin{frame}
    \frametitle{Циклічні групи}
    \begin{block}{Означення}
        Групою $G$ називається циклічною групою, якщо:
       \begin{equation*}
           \exists g \in G, \forall a\in G, \exists k\in Z: a=g^k
       \end{equation*}
       Елементом $g$ називають твірним (породжуючим) елементом групи $G$, або її генератором,а групу $G$ -- циклічною групою, породженою елементом $g$, і це позначають як $G=<g>$.
    \end{block}
    \begin{figure}
        \centering
        \includegraphics[width=1\linewidth]{2-crypto-sym/images/cells_small.png}
        \label{fig:placeholder}
    \end{figure}
\end{frame}
%------------------------------------------------
\begin{frame}
    \frametitle{Теорія кілець}
    \begin{block}{Означення}
        \textbf{Кільцем} називається алгебраїчна структура з двома операціями, які умовно назвемо додаванням та множенням. При цьому:
        \begin{enumerate}
            \item відносно операції додавання $R$ утворює абелеву групу; (нейтральний елемент цієї групи позначається $0$, обернений до $a$ через $-a$, він називається протилежним до   елементом);
            \item відносно операції множення $R$ утворює напівгрупу;
            \item операції додавання та множення пов’язані законами дистрибутивності: 
            \begin{align*}
                (a+b)\cdot c=a\cdot c+b\cdot c \\
                c\cdot (a+b)=c\cdot a+c\cdot b
            \end{align*}
        \end{enumerate}
         Кількість елементів кільця називають порядком кільця.
    \end{block}
\end{frame}
%------------------------------------------------
\begin{frame}
    \frametitle{Теорія полів}
    \begin{block}{Означення}
        \textbf{Полем} називається алгебраїчна структура з двома операціями: додаванням і множенням, що задовольняють умовам:
        \begin{enumerate}
            \item за додаванням $F$ є абелевою групою;
            \item за множенням усі ненульові елементи $F$ також утворюють абелеву групу; 
            \item 3) операції додавання та множення пов’язані законами дистрибутивності: 
            \begin{align*}
                (a+b)\cdot c=a\cdot c+b\cdot c \\
                c\cdot (a+b)=c\cdot a+c\cdot b
            \end{align*}
        \end{enumerate}
        З означення поля видно, що на відміну від комутативного кільця з одиницею, у полі кожен ненульовий елемент є оборотним.
    \end{block}
\end{frame}
%------------------------------------------------
\begin{frame}
    \frametitle{Теорія полів}
    \begin{block}{Означення}
        Скінченні поля називають полями Галуа (Galois Fields) і позначають $GF(p^n)$ або $F_{p^n}$, де $p^n$ -- кількість елементів поля.

        Тепер замість $Z_p$ будемо використовувати позначення $F_p$.
    \end{block}
    \begin{figure}
        \centering
        \includegraphics[width=1\linewidth]{2-crypto-sym/images/cells_small.png}
    \end{figure}
\end{frame}
\section{Асиметрична криптографія}
%------------------------------------------------
\begin{frame}
    \frametitle{Задачі}
        \begin{block}{}
	\textbf{1.} Розглянемо алгоритм пошуку генератора мультиплікативної групи скінченного поля $F_q^*$:
    \begin{enumerate}
        \item обрати випадковий елемент $g\in_{R}F_q^*$;
        \item перевірити чи є цей елемент генератором;
        \item якщо так, повернути $g$; інакше, повернутись на перший крок.
    \end{enumerate}
    Скільки елементів в середньому буде перевіряти даний алгоритм, якщо елементи обираються із поверненням? Дати відповідь в загальному випадку та при $q=19,27,256$.
    	\end{block}
            \begin{figure}
        \centering
        \includegraphics[width=1\linewidth]{2-crypto-sym/images/cells_small.png}
    \end{figure}
\end{frame}
%------------------------------------------------
\begin{frame}
    \frametitle{Задачі}
        \begin{block}{}
        \textbf{2.} Розглянемо алгоритм випадкового пошуку генератора $F_q^*$ з попередньої задачі, але елементи для перевірки обираються без повернення. Скільки ітерацій зробить в середньому такий алгоритм при пошуку генераторів при  $q=7,8$?
    	\end{block}
                \begin{figure}
        \centering
        \includegraphics[width=1\linewidth]{2-crypto-sym/images/cells_small.png}
    \end{figure}
\end{frame}
%------------------------------------------------
\begin{frame}
    \frametitle{Задачі}
        \begin{block}{}
        \textbf{3.} Обчислити $(\frac{19}{51}),(\frac{52}{41}),(\frac{616}{3315})$
        \end{block}
        \begin{figure}
            \centering
            \includegraphics[width=1\linewidth]{2-crypto-sym/images/cells.png}
        \end{figure}
\end{frame}
%------------------------------------------------
\begin{frame}
    \frametitle{Задачі}
        \begin{block}{}
        \textbf{4.}	(Критерій Ойлера) Нехай  $p$ -- просте, $a\in Z$. Доведіть, що 
        \begin{align*}
            (\frac{a}{p})\equiv a^{\frac{p-1}{2}}(\text{mod }p)
        \end{align*}
    	\end{block}
        \begin{figure}
        \centering
        \includegraphics[width=1\linewidth]{2-crypto-sym/images/cells_small.png}
    \end{figure}
\end{frame}
%------------------------------------------------
\begin{frame}
    \frametitle{Задачі}
        \begin{block}{}
        \textbf{5.} Для деякого цілого $k$ доведіть, що 
        \begin{enumerate}
            \item $-1$ є квадратичним нелишком за простим модулем $р$ тоді та тільки тоді, коли $p=4k+3$ 
            \item $2$ є квадратичним нелишком за простим модулем $р$ при $p=8k+5$. 
        \end{enumerate}
    	\end{block}    
            \begin{figure}
        \centering
        \includegraphics[width=1\linewidth]{2-crypto-sym/images/cells_small.png}
    \end{figure}
\end{frame}
%------------------------------------------------
\begin{frame}
    \frametitle{Задачі}
        \begin{block}{}
        \textbf{6.} Нехай $р$ -- просте, $a\in Q_p$. Скільки квадратичних лишків та квадратичних нелишків буде міститись серед квадратних коренів $а$? Дати відповідь у випадку, коли $р$ має вид:
        \begin{enumerate}
            \item $p=4k+3$
            \item $p=4k+1$
        \end{enumerate}
        \begin{figure}
            \centering
            \includegraphics[width=0.8\linewidth]{2-crypto-sym/images/cells.png}
        \end{figure}

    \end{block}
\end{frame}
%------------------------------------------------

    \begin{frame}[plain, standout]
    \begin{figure}
        \centering
        \includegraphics[width=0.9\linewidth]{Cryptocats.jpg}
        \label{fig:placeholder}
    \end{figure}
    \centering
    \LARGE
    \textbf{Питання?} \\

    \vspace{0.2cm} \Huge \ding{170} \large \\

  \end{frame}
\end{document}